\[ %%%%%%%%%%%%%%%%%%%%%%%%%%%%%%%%%%%%%%%%%%%%%%%%%%%%%%%%%%%%%%%%%%%%%%%%%%%%%%%% \subsection{Long-Horizon Forecasting of Nonlinear Dynamical Systems} \label{sec:long_horizon} %%%%%%%%%%%%%%%%%%%%%%%%%%%%%%%%%%%%%%%%%%%%%%%%%%%%%%%%%%%%%%%%%%%%%%%%%%%%%%%% Long-horizon forecasting of nonlinear dynamical systems represents one of the most challenging problems in scientific machine learning. While modern neural architectures have achieved remarkable accuracy for short-term prediction tasks, maintaining physically consistent and stable trajectories over extended temporal horizons remains a fundamental unresolved challenge. Many scientific systems of practical importance exhibit nonlinear behavior characterized by sensitivity to initial conditions, multi-scale interactions, and complex temporal dependencies. Examples include turbulent fluid flows, atmospheric and climate dynamics, plasma evolution, biological systems, chemical reaction networks, and engineered infrastructures. In such systems, small local approximation errors introduced during prediction may propagate through nonlinear interactions and produce significant deviations from the true trajectory over time. A common strategy for learning dynamical systems is autoregressive forecasting. Given a system state $x_t$, a learned evolution operator predicts the subsequent state: \begin{equation} x_{t+1}=F_\theta(x_t), \end{equation} where $F_\theta$ represents a neural approximation of the underlying dynamical evolution operator. Long-term prediction is then obtained by recursively applying the learned operator: \begin{equation} x_{t+k}=F_\theta^{k}(x_t). \end{equation} Although this formulation is computationally efficient, repeated application of an approximate evolution operator introduces cumulative prediction errors. Even when the one-step prediction error is small, deviations may amplify through repeated iterations, especially in systems with chaotic or unstable dynamics. The mathematical origin of this phenomenon is closely related to the sensitivity properties of nonlinear dynamical systems. For systems characterized by positive Lyapunov exponents, infinitesimal perturbations may grow exponentially according to: \begin{equation} ||\delta x(t)|| \approx ||\delta x(0)||e^{\lambda t}, \end{equation} where $\lambda$ denotes the dominant Lyapunov exponent. Consequently, a machine learning model that does not preserve the intrinsic stability properties of the underlying dynamics may gradually diverge from physically meaningful trajectories. Several approaches have been proposed to address long-horizon forecasting challenges. Sequence-based models, including recurrent neural networks and long short-term memory architectures, demonstrated the ability to capture temporal dependencies in complex systems. More recently, transformer-based architectures have achieved significant success in large-scale time-series prediction due to their ability to model long-range dependencies through attention mechanisms. However, purely data-driven temporal models often face difficulties when extrapolating beyond the statistical distribution represented in training datasets. Since their predictions are optimized primarily according to empirical data consistency, they may generate trajectories that appear statistically plausible while violating fundamental physical properties such as conservation laws, stability constraints, or dynamical invariants. Reservoir computing approaches, including echo state networks, have also been investigated for forecasting chaotic systems. These methods demonstrated promising results for certain low-dimensional nonlinear systems by exploiting high-dimensional latent representations. Nevertheless, their performance often deteriorates for large-scale spatially distributed systems where explicit representation of physical interactions becomes necessary. Recent developments in scientific machine learning have attempted to improve long-horizon forecasting through physics-informed constraints, neural operators, and hybrid numerical-learning architectures. These approaches have demonstrated improved stability compared with purely data-driven methods; however, many existing frameworks still rely on training-time constraints and provide limited mechanisms for actively controlling error propagation during inference. A fundamental challenge therefore remains: a reliable long-horizon scientific predictor must not only approximate the instantaneous evolution of a system but also preserve the structural properties that govern its future trajectory. This requires computational architectures capable of maintaining physical consistency, regulating uncertainty growth, and adapting information propagation throughout the entire forecasting process. The PHYSGEN framework addresses this challenge by introducing a physics-guided graph-based approach in which long-term prediction stability is treated as an architectural objective rather than an emergent property of training. By combining graph-based representation of physical interactions, multi-level physics guidance, and adaptive stability mechanisms, PHYSGEN aims to provide a computational foundation for robust long-horizon prediction of nonlinear dynamical systems. \] 
